\documentclass[11pt]{article}

\usepackage[utf8]{inputenc}
\usepackage[T1]{fontenc}
\usepackage[spanish,mexico]{babel}
\usepackage{amsmath,amssymb}
\usepackage{geometry}
\usepackage{enumitem}
\usepackage{multicol}
\usepackage{fancyhdr}

\geometry{letterpaper,margin=1.7cm}
\setlength{\parindent}{0pt}

\pagestyle{fancy}
\fancyhf{}
\lhead{Álgebra Lineal}
\rhead{Fluidez y Agilidad Aritmética --- Versión A}
\cfoot{\thepage}

\begin{document}

\begin{center}
{\Large\bfseries Prueba de Fluidez y Agilidad Aritmética}\\[1mm]
{\large Versión A --- Diagnóstico preliminar}
\end{center}

\vspace{1mm}

\textbf{Nombre:} \rule{9cm}{0.4pt}\hfill
\textbf{Fecha:} \rule{3cm}{0.4pt}

\vspace{2mm}

\textbf{Instrucciones.}
Resuelva mentalmente cada operación y escriba únicamente el resultado.
No utilice calculadora. Trabaje con rapidez, pero procure mantener la
precisión. Si una respuesta es una fracción, simplifíquela.

\medskip

\textbf{Propósito diagnóstico.}
El instrumento explora fluidez aritmética general mediante operaciones
elementales, jerarquía de operaciones, signos, potencias, fracciones,
decimales, porcentajes y operaciones mixtas. No constituye por sí solo
una prueba psicométrica estandarizada.

\medskip

\textbf{Tiempo de aplicación:} \rule{2.5cm}{0.4pt}
\hfill
\textbf{Reactivos contestados:} \rule{1.5cm}{0.4pt}
\hfill
\textbf{Aciertos:} \rule{1.5cm}{0.4pt}

\subsection*{I. Aritmética inmediata}
\begin{multicols}{2}\begin{enumerate}[leftmargin=*,start=1]
\item $\displaystyle 7+6=$ \rule{2.2cm}{0.4pt}
\item $\displaystyle 8+9=$ \rule{2.2cm}{0.4pt}
\item $\displaystyle 15-8=$ \rule{2.2cm}{0.4pt}
\item $\displaystyle 13-7=$ \rule{2.2cm}{0.4pt}
\item $\displaystyle 6\cdot4=$ \rule{2.2cm}{0.4pt}
\item $\displaystyle 7\cdot8=$ \rule{2.2cm}{0.4pt}
\item $\displaystyle 36\div6=$ \rule{2.2cm}{0.4pt}
\item $\displaystyle 42\div7=$ \rule{2.2cm}{0.4pt}
\item $\displaystyle 12+9=$ \rule{2.2cm}{0.4pt}
\item $\displaystyle 17-9=$ \rule{2.2cm}{0.4pt}
\item $\displaystyle 5\cdot9=$ \rule{2.2cm}{0.4pt}
\item $\displaystyle 48\div8=$ \rule{2.2cm}{0.4pt}
\item $\displaystyle 14+17=$ \rule{2.2cm}{0.4pt}
\item $\displaystyle 25-16=$ \rule{2.2cm}{0.4pt}
\item $\displaystyle 9\cdot6=$ \rule{2.2cm}{0.4pt}
\item $\displaystyle 63\div9=$ \rule{2.2cm}{0.4pt}
\item $\displaystyle 18+7-5=$ \rule{2.2cm}{0.4pt}
\item $\displaystyle 20-6+3=$ \rule{2.2cm}{0.4pt}
\item $\displaystyle 4+7+8=$ \rule{2.2cm}{0.4pt}
\item $\displaystyle 30-8-7=$ \rule{2.2cm}{0.4pt}
\end{enumerate}\end{multicols}
\subsection*{II. Operaciones combinadas}
\begin{multicols}{2}\begin{enumerate}[leftmargin=*,start=21]
\item $\displaystyle 3\cdot5+4=$ \rule{2.2cm}{0.4pt}
\item $\displaystyle 4\cdot6-7=$ \rule{2.2cm}{0.4pt}
\item $\displaystyle 18\div3+5=$ \rule{2.2cm}{0.4pt}
\item $\displaystyle 24\div6-2=$ \rule{2.2cm}{0.4pt}
\item $\displaystyle 5+3\cdot4=$ \rule{2.2cm}{0.4pt}
\item $\displaystyle 20-2\cdot6=$ \rule{2.2cm}{0.4pt}
\item $\displaystyle 2(5+3)=$ \rule{2.2cm}{0.4pt}
\item $\displaystyle 3(7-4)=$ \rule{2.2cm}{0.4pt}
\item $\displaystyle 18\div(6-3)=$ \rule{2.2cm}{0.4pt}
\item $\displaystyle 4(6)-3(5)=$ \rule{2.2cm}{0.4pt}
\end{enumerate}\end{multicols}
\subsection*{III. Signos}
\begin{multicols}{2}\begin{enumerate}[leftmargin=*,start=31]
\item $\displaystyle -3+8=$ \rule{2.2cm}{0.4pt}
\item $\displaystyle 5-9=$ \rule{2.2cm}{0.4pt}
\item $\displaystyle -4-3=$ \rule{2.2cm}{0.4pt}
\item $\displaystyle -7+2=$ \rule{2.2cm}{0.4pt}
\item $\displaystyle 6-(-3)=$ \rule{2.2cm}{0.4pt}
\item $\displaystyle -5-(-2)=$ \rule{2.2cm}{0.4pt}
\item $\displaystyle (-3)(4)=$ \rule{2.2cm}{0.4pt}
\item $\displaystyle (-5)(-2)=$ \rule{2.2cm}{0.4pt}
\item $\displaystyle 24\div(-6)=$ \rule{2.2cm}{0.4pt}
\item $\displaystyle (-18)\div(-3)=$ \rule{2.2cm}{0.4pt}
\end{enumerate}\end{multicols}
\subsection*{IV. Potencias y jerarquía de operaciones}
\begin{multicols}{2}\begin{enumerate}[leftmargin=*,start=41]
\item $\displaystyle 3^2+4=$ \rule{2.2cm}{0.4pt}
\item $\displaystyle 2^3-3=$ \rule{2.2cm}{0.4pt}
\item $\displaystyle 5^2-20=$ \rule{2.2cm}{0.4pt}
\item $\displaystyle 2^4\div4=$ \rule{2.2cm}{0.4pt}
\item $\displaystyle 3(2^2)-5=$ \rule{2.2cm}{0.4pt}
\item $\displaystyle 2(3^2)-4=$ \rule{2.2cm}{0.4pt}
\item $\displaystyle 4+2(5-3)=$ \rule{2.2cm}{0.4pt}
\item $\displaystyle 3(6-2)-5=$ \rule{2.2cm}{0.4pt}
\item $\displaystyle 2^3+3^2=$ \rule{2.2cm}{0.4pt}
\item $\displaystyle 5^2-3^2=$ \rule{2.2cm}{0.4pt}
\end{enumerate}\end{multicols}
\subsection*{V. Fracciones: suma y resta}
\begin{multicols}{2}\begin{enumerate}[leftmargin=*,start=51]
\item $\displaystyle \frac12+\frac14=$ \rule{2.2cm}{0.4pt}
\item $\displaystyle \frac34-\frac12=$ \rule{2.2cm}{0.4pt}
\item $\displaystyle \frac13+\frac13=$ \rule{2.2cm}{0.4pt}
\item $\displaystyle \frac23-\frac13=$ \rule{2.2cm}{0.4pt}
\item $\displaystyle \frac12+\frac13=$ \rule{2.2cm}{0.4pt}
\item $\displaystyle \frac34+\frac14=$ \rule{2.2cm}{0.4pt}
\item $\displaystyle \frac56-\frac13=$ \rule{2.2cm}{0.4pt}
\item $\displaystyle \frac23+\frac16=$ \rule{2.2cm}{0.4pt}
\item $\displaystyle \frac78-\frac38=$ \rule{2.2cm}{0.4pt}
\item $\displaystyle 1-\frac14=$ \rule{2.2cm}{0.4pt}
\item $\displaystyle 1-\frac38=$ \rule{2.2cm}{0.4pt}
\item $\displaystyle 2-\frac34=$ \rule{2.2cm}{0.4pt}
\item $\displaystyle \frac23+\frac14=$ \rule{2.2cm}{0.4pt}
\item $\displaystyle \frac56-\frac14=$ \rule{2.2cm}{0.4pt}
\item $\displaystyle \frac35+\frac12=$ \rule{2.2cm}{0.4pt}
\end{enumerate}\end{multicols}
\subsection*{VI. Fracciones: productos, cocientes y operaciones combinadas}
\begin{multicols}{2}\begin{enumerate}[leftmargin=*,start=66]
\item $\displaystyle \frac12\cdot\frac34=$ \rule{2.2cm}{0.4pt}
\item $\displaystyle \frac23\cdot\frac32=$ \rule{2.2cm}{0.4pt}
\item $\displaystyle \frac35\cdot\frac52=$ \rule{2.2cm}{0.4pt}
\item $\displaystyle \frac34\div\frac12=$ \rule{2.2cm}{0.4pt}
\item $\displaystyle \frac23\div\frac43=$ \rule{2.2cm}{0.4pt}
\item $\displaystyle \frac12+\frac13+\frac16=$ \rule{2.2cm}{0.4pt}
\item $\displaystyle \frac34-\frac12+\frac14=$ \rule{2.2cm}{0.4pt}
\item $\displaystyle 1-\frac12+\frac14=$ \rule{2.2cm}{0.4pt}
\item $\displaystyle 2-\frac12-\frac14=$ \rule{2.2cm}{0.4pt}
\item $\displaystyle \frac23+\frac12-\frac16=$ \rule{2.2cm}{0.4pt}
\item $\displaystyle \frac12\left(\frac34\right)=$ \rule{2.2cm}{0.4pt}
\item $\displaystyle 2\left(\frac34\right)=$ \rule{2.2cm}{0.4pt}
\item $\displaystyle 3\left(\frac23\right)=$ \rule{2.2cm}{0.4pt}
\item $\displaystyle \frac12\div\frac14=$ \rule{2.2cm}{0.4pt}
\item $\displaystyle \frac34\div\frac38=$ \rule{2.2cm}{0.4pt}
\end{enumerate}\end{multicols}
\subsection*{VII. Decimales y porcentajes}
\begin{multicols}{2}\begin{enumerate}[leftmargin=*,start=81]
\item $\displaystyle 0.5+0.25=$ \rule{2.2cm}{0.4pt}
\item $\displaystyle 1.2+0.8=$ \rule{2.2cm}{0.4pt}
\item $\displaystyle 2.5-1.5=$ \rule{2.2cm}{0.4pt}
\item $\displaystyle 0.4(5)=$ \rule{2.2cm}{0.4pt}
\item $\displaystyle 0.5(18)=$ \rule{2.2cm}{0.4pt}
\item $\displaystyle 25\%\text{ de }20=$ \rule{2.2cm}{0.4pt}
\item $\displaystyle 50\%\text{ de }18=$ \rule{2.2cm}{0.4pt}
\item $\displaystyle 10\%\text{ de }70=$ \rule{2.2cm}{0.4pt}
\item $\displaystyle 20\%\text{ de }30=$ \rule{2.2cm}{0.4pt}
\item $\displaystyle 75\%\text{ de }8=$ \rule{2.2cm}{0.4pt}
\end{enumerate}\end{multicols}
\subsection*{VIII. Operaciones mixtas}
\begin{multicols}{2}\begin{enumerate}[leftmargin=*,start=91]
\item $\displaystyle \frac12+0.25=$ \rule{2.2cm}{0.4pt}
\item $\displaystyle \frac34-0.5=$ \rule{2.2cm}{0.4pt}
\item $\displaystyle 50\%\text{ de }\frac32=$ \rule{2.2cm}{0.4pt}
\item $\displaystyle 2-\frac12(3)=$ \rule{2.2cm}{0.4pt}
\item $\displaystyle \frac34+\frac12-\frac14=$ \rule{2.2cm}{0.4pt}
\item $\displaystyle \left(\frac12+\frac14\right)\div\frac34=$ \rule{2.2cm}{0.4pt}
\item $\displaystyle \frac23\left(1+\frac12\right)=$ \rule{2.2cm}{0.4pt}
\item $\displaystyle 1-\left(\frac12-\frac14\right)=$ \rule{2.2cm}{0.4pt}
\item $\displaystyle \frac12+\frac12\div\frac14=$ \rule{2.2cm}{0.4pt}
\item $\displaystyle \left(1-\frac13\right)\div\frac23=$ \rule{2.2cm}{0.4pt}
\end{enumerate}\end{multicols}

\newpage

\section*{Clave de respuestas --- uso del profesor}

\subsection*{I. Aritmética inmediata}
\begin{multicols}{4}\begin{enumerate}[leftmargin=*,start=1]
\item $\displaystyle 13$
\item $\displaystyle 17$
\item $\displaystyle 7$
\item $\displaystyle 6$
\item $\displaystyle 24$
\item $\displaystyle 56$
\item $\displaystyle 6$
\item $\displaystyle 6$
\item $\displaystyle 21$
\item $\displaystyle 8$
\item $\displaystyle 45$
\item $\displaystyle 6$
\item $\displaystyle 31$
\item $\displaystyle 9$
\item $\displaystyle 54$
\item $\displaystyle 7$
\item $\displaystyle 20$
\item $\displaystyle 17$
\item $\displaystyle 19$
\item $\displaystyle 15$
\end{enumerate}\end{multicols}
\subsection*{II. Operaciones combinadas}
\begin{multicols}{4}\begin{enumerate}[leftmargin=*,start=21]
\item $\displaystyle 19$
\item $\displaystyle 17$
\item $\displaystyle 11$
\item $\displaystyle 2$
\item $\displaystyle 17$
\item $\displaystyle 8$
\item $\displaystyle 16$
\item $\displaystyle 9$
\item $\displaystyle 6$
\item $\displaystyle 9$
\end{enumerate}\end{multicols}
\subsection*{III. Signos}
\begin{multicols}{4}\begin{enumerate}[leftmargin=*,start=31]
\item $\displaystyle 5$
\item $\displaystyle -4$
\item $\displaystyle -7$
\item $\displaystyle -5$
\item $\displaystyle 9$
\item $\displaystyle -3$
\item $\displaystyle -12$
\item $\displaystyle 10$
\item $\displaystyle -4$
\item $\displaystyle 6$
\end{enumerate}\end{multicols}
\subsection*{IV. Potencias y jerarquía de operaciones}
\begin{multicols}{4}\begin{enumerate}[leftmargin=*,start=41]
\item $\displaystyle 13$
\item $\displaystyle 5$
\item $\displaystyle 5$
\item $\displaystyle 4$
\item $\displaystyle 7$
\item $\displaystyle 14$
\item $\displaystyle 8$
\item $\displaystyle 7$
\item $\displaystyle 17$
\item $\displaystyle 16$
\end{enumerate}\end{multicols}
\subsection*{V. Fracciones: suma y resta}
\begin{multicols}{4}\begin{enumerate}[leftmargin=*,start=51]
\item $\displaystyle \frac34$
\item $\displaystyle \frac14$
\item $\displaystyle \frac23$
\item $\displaystyle \frac13$
\item $\displaystyle \frac56$
\item $\displaystyle 1$
\item $\displaystyle \frac12$
\item $\displaystyle \frac56$
\item $\displaystyle \frac12$
\item $\displaystyle \frac34$
\item $\displaystyle \frac58$
\item $\displaystyle \frac54$
\item $\displaystyle \frac{11}{12}$
\item $\displaystyle \frac7{12}$
\item $\displaystyle \frac{11}{10}$
\end{enumerate}\end{multicols}
\subsection*{VI. Fracciones: productos, cocientes y operaciones combinadas}
\begin{multicols}{4}\begin{enumerate}[leftmargin=*,start=66]
\item $\displaystyle \frac38$
\item $\displaystyle 1$
\item $\displaystyle 1$
\item $\displaystyle \frac32$
\item $\displaystyle \frac12$
\item $\displaystyle 1$
\item $\displaystyle \frac12$
\item $\displaystyle \frac34$
\item $\displaystyle \frac54$
\item $\displaystyle 1$
\item $\displaystyle \frac38$
\item $\displaystyle \frac32$
\item $\displaystyle 2$
\item $\displaystyle 2$
\item $\displaystyle 2$
\end{enumerate}\end{multicols}
\subsection*{VII. Decimales y porcentajes}
\begin{multicols}{4}\begin{enumerate}[leftmargin=*,start=81]
\item $\displaystyle 0.75$
\item $\displaystyle 2$
\item $\displaystyle 1$
\item $\displaystyle 2$
\item $\displaystyle 9$
\item $\displaystyle 5$
\item $\displaystyle 9$
\item $\displaystyle 7$
\item $\displaystyle 6$
\item $\displaystyle 6$
\end{enumerate}\end{multicols}
\subsection*{VIII. Operaciones mixtas}
\begin{multicols}{4}\begin{enumerate}[leftmargin=*,start=91]
\item $\displaystyle 0.75$
\item $\displaystyle 0.25$
\item $\displaystyle 0.75$
\item $\displaystyle \frac12$
\item $\displaystyle 1$
\item $\displaystyle 1$
\item $\displaystyle 1$
\item $\displaystyle \frac34$
\item $\displaystyle \frac52$
\item $\displaystyle 1$
\end{enumerate}\end{multicols}

\end{document}
